\section*{अध्याय \ref{ch:g5:data} --- समस्याएँ और आलेख}

\begin{solution}{exo:g5:data:1}
टिकट: $32 \times 7 = 224$। कुल: $240 + 224 = 464$।
\end{solution}

\begin{solution}{exo:g5:data:2}
एक kg का: $7.50 \div 3 = 2.50$। पाँच kg:
$5 \times 2.50 = 12.50$।
\end{solution}

\begin{solution}{exo:g5:data:3}
$87 \div 6 = 14.5$: हर एक $14.50$ देता है।
\end{solution}

\begin{solution}{exo:g5:data:4}
ख़र्च: $13.90 + 8.50 = 22.40$। बचे: $25 - 22.40 = 2.60$।
\end{solution}

\begin{solution}{exo:g5:data:5}
$50$ mm से ज़्यादा: जनवरी ($60$), मार्च ($55$), अप्रैल ($70$),
मई ($80$)। मई में से जून: $80 - 30 = 50$ mm।
\end{solution}

\begin{solution}{exo:g5:data:6}
हफ़्ता $2$: $6$ cm। $4$ cm की बढ़त: हफ़्ते $2$ और $3$ के बीच
($6 \to 10$), और हफ़्ते $4$ और $5$ के बीच ($15 \to 19$)।
\end{solution}

\begin{solution}{exo:g5:data:7}
$4$, $7$, $3$, $6$ ऊँचाइयों की चार पट्टियाँ (अंकन हर $1$ या
$2$ गोल पर)। कुल: $4 + 7 + 3 + 6 = 20$ गोल।
\end{solution}

\begin{solution}{exo:g5:data:8}
पैकेट में एक का: $3.20 \div 8 = 0.40$। हर दही पर बचत:
$0.50 - 0.40 = 0.10$ (दस सेंट)।
\end{solution}

\begin{solution}{exo:g5:data:9}
पहला तरीक़ा: शनिवार $148 \times 9.50 = 1\,406$; रविवार
$95 \times 9.50 = 902.50$; \omterm{ex:g1:subtraction:difference}{अंतर} $503.50$।

दूसरा तरीक़ा: $148 - 95 = 53$ टिकट ज़्यादा, हर एक $9.50$ का:
$53 \times 9.50 = 503.50$। वही उत्तर --- कमाइयों का \omterm{ex:g1:subtraction:difference}{अंतर},
\omterm{ex:g1:subtraction:difference}{अंतर} की कमाई ही है।
\end{solution}

\begin{solution}{exo:g5:data:10}
एक व्यक्ति के लिए: $300 \div 4 = 75$ g। दस के लिए:
$10 \times 75 = 750$ g चावल।
\end{solution}

\begin{solution}{exo:g5:data:11}
कई सही समस्याएँ बन सकती हैं --- जैसे: ``मिलो $1.20$-$1.20$ के
$3$ जूस और $3.40$ का एक सैंडविच ख़रीदता है। वह $12.50$ का एक
कूपन देता है। कूपन में कितना बचा?'' हल: जूस
$3 \times 1.20 = 3.60$; कुल ख़र्च $3.60 + 3.40 = 7$; बचा
$12.50 - 7 = 5.50$।

\end{solution}
